\chapter{Descripción del problema}
\label{cap:problema}

Este capítulo delimita el problema que aborda el trabajo. Primero se formula
el problema y se fija el patrón objetivo de transformación; después se define
qué se entiende exactamente por condicional anidada \emph{combinable} y se
muestran casos próximos que no lo son; por último, se sitúa el patrón dentro
de una familia más amplia de problemas y se acota el alcance del prototipo.
La metodología, la operacionalización de las preguntas de investigación y las
hipótesis de trabajo se desarrollan en el \autoref{cap:metodologia}.

\section{Formulación del problema}
\label{sec:formulacion}

Los condicionales anidados son una de las construcciones que más contribuyen
a la complejidad cognitiva de un método: cada nivel de anidamiento añade un
salto mental que el lector debe sostener~\cite{10.1145/3194164.3194186}
(\autoref{sec:cc-antecedentes}). Sin embargo, en muchos casos dos sentencias
\lstinline|if| anidadas sin ramas alternativas no expresan dos decisiones
distintas, sino una sola condición conjuntiva escrita en dos niveles. Cuando
esto ocurre, aplanarlas en \lstinline|if (A && B)| simplifica el código sin
cambiar lo que hace. La misma idea se extiende a cadenas de más de dos
\lstinline|if| anidados, que se combinan aplicando la regla por pares de forma
sucesiva (\autoref{sec:alcance}).

El problema abordado nace de esa observación y tiene dos dimensiones. La
primera consiste en determinar en qué condiciones esa combinación es
semánticamente segura y automatizable. La segunda consiste en medir su efecto
sobre la complejidad cognitiva. Interesa responder a ambas no solo para una
herramienta determinista, sino también para un gran modelo de lenguaje, lo
que permite comparar dos formas muy distintas de afrontar la misma tarea.
Estas dos dimensiones se recogen en las preguntas de investigación del
trabajo, enunciadas en la \autoref{sec:rq-intro}.

\section{Patrón objetivo}
\label{sec:patron}

La transformación canónica estudiada parte de un anidamiento simple y se
transforma en una condición conjuntiva equivalente, como ilustra la
\autoref{fig:before-after}.

\begin{figure}[ht]
\centering
\begin{minipage}[c]{0.44\linewidth}
\centering
\textbf{Antes}\par\smallskip
\begin{lstlisting}[numbers=none]
if (A) {
    if (B) {
        S
    }
}
\end{lstlisting}
\end{minipage}%
\begin{minipage}[c]{0.10\linewidth}
\centering$\Longrightarrow$
\end{minipage}%
\begin{minipage}[c]{0.44\linewidth}
\centering
\textbf{Después}\par\smallskip
\begin{lstlisting}[numbers=none]
if (A && B) {
    S
}
\end{lstlisting}
\end{minipage}
\caption{Esquema \emph{before/after} del patrón objetivo: dos \texttt{if}
anidados (izquierda) se combinan en una única condición conjuntiva
\lstinline|A && B| (derecha). La transformación preserva el orden de
evaluación y es semánticamente equivalente bajo las precondiciones P1--P5.}
\label{fig:before-after}
\end{figure}

\noindent donde \lstinline|A| y \lstinline|B| son expresiones booleanas y
\lstinline|S| un bloque de sentencias arbitrario. La equivalencia se sostiene
bajo las precondiciones de aplicabilidad (\autoref{sec:precondiciones}): el
operador \lstinline|&&| evalúa primero \lstinline|A| y, solo si se cumple,
\lstinline|B|, preservando así el orden y el cortocircuito del anidamiento
original y, con ello, el comportamiento observable.

Las condiciones \lstinline|A| y \lstinline|B| pueden ser expresiones booleanas
compuestas. En ese caso, la combinación preserva la precedencia original
añadiendo paréntesis cuando es necesario (por ejemplo, si una de ellas es una
disyunción que pasa a formar parte de la conjunción).

Para hacer tangible el efecto sobre la métrica, conviene contar la complejidad
cognitiva de ambas versiones según el modelo de SonarSource
(\autoref{sec:cc-antecedentes}). En la versión anidada, el \lstinline|if|
externo aporta un incremento estructural (+1) y el \lstinline|if| interno
aporta otro estructural más uno de anidamiento por estar un nivel más adentro
(+2): en total, 3. En la versión combinada, el único \lstinline|if| aporta +1 y
la secuencia de operadores lógicos \lstinline|&&| aporta +1: en total, 2. La
combinación ahorra, por tanto, el incremento de anidamiento del segundo
\lstinline|if|. La \autoref{fig:cc-desglose} ilustra este desglose sobre un
método concreto, anotando la contribución de cada nodo a la complejidad
cognitiva del método.

\begin{figure}[ht]
\centering
\begin{minipage}{0.9\linewidth}
\textbf{Antes (anidado): CC del método $= 3$}
\begin{lstlisting}[numbers=none]
boolean puedeEnviar(boolean pagado, boolean enStock) {
    if (pagado) {            // +1  (incremento estructural)
        if (enStock) {       // +2  (estructural +1, anidamiento +1)
            return true;
        }
    }
    return false;
}
\end{lstlisting}
\smallskip
\textbf{Después (combinado): CC del método $= 2$}
\begin{lstlisting}[numbers=none]
boolean puedeEnviar(boolean pagado, boolean enStock) {
    if (pagado && enStock) { // +1 (if) y +1 (secuencia &&)
        return true;
    }
    return false;
}
\end{lstlisting}
\end{minipage}
\caption{Desglose de la complejidad cognitiva por nodo según el modelo de
SonarSource. En la versión anidada, el \texttt{if} externo aporta $+1$ y el
\texttt{if} interno $+2$ (incremento estructural más uno de anidamiento por
estar un nivel más adentro), de modo que el método suma CC${}=3$; al combinar
las condiciones, el único \texttt{if} aporta $+1$ y la conjunción
\lstinline|&&| otro $+1$, con CC${}=2$. La combinación elimina el incremento de
anidamiento del segundo \texttt{if}.}
\label{fig:cc-desglose}
\end{figure}

\section{Qué es una condicional anidada combinable}
\label{sec:definicion-combinable}

Se denomina \emph{condicional anidada combinable} a un par de sentencias
\lstinline|if| en el que el \lstinline|if| externo no tiene rama alternativa,
su rama \lstinline|then| contiene exactamente otra sentencia \lstinline|if|, y
esta segunda tampoco tiene rama alternativa. Para un par así, sustituir el
anidamiento por \lstinline|if (A && B)| preserva el comportamiento observable,
porque \lstinline|&&| evalúa \lstinline|A| y, solo si es cierto,
\lstinline|B|: el mismo orden y el mismo cortocircuito que en la versión
anidada. Esta es la condición de combinabilidad \emph{estructural}.

El prototipo añade una restricción adicional: exige que las condiciones estén
libres de efectos colaterales. No se debe a que la combinación deje de
preservar el comportamiento en su presencia, sino a que solo bajo esa restricción puede garantizarse la
seguridad de forma automática y verificable sobre la estructura del código,
sin recurrir a un análisis semántico que excede el alcance del prototipo. El
conjunto completo de precondiciones se formaliza en la
\autoref{sec:precondiciones}.

\section{Casos que no son combinables}
\label{sec:no-combinables}

La dificultad del problema no está en el caso favorable, sino en distinguirlo
de otros muy parecidos en los que aplanar el anidamiento \emph{cambia el
comportamiento}. Dos ejemplos permiten verlo con claridad.

Si el \lstinline|if| externo tiene una rama \lstinline|else|, esta cubre el
caso en que \lstinline|A| es falso; combinar las condiciones la haría
aplicable también cuando \lstinline|A| es cierto pero \lstinline|B| es falso,
alterando el flujo:

\begin{lstlisting}[caption={No combinable: \texttt{else} en el \texttt{if} externo.},label={lst:no-else}]
if (A) {
    if (B) { S }
} else {
    T          // se ejecuta cuando !A; "if (A && B) ... else T" lo rompe
}
\end{lstlisting}

De forma análoga, si el bloque externo contiene alguna sentencia además del
\lstinline|if| interno, esa sentencia se ejecuta siempre que se cumpla
\lstinline|A|; tras la combinación pasaría a ejecutarse solo cuando además se
cumpla \lstinline|B|:

\begin{lstlisting}[caption={No combinable: sentencia adicional en el bloque externo.},label={lst:no-extra}]
if (A) {
    T;             // se ejecuta cuando A, con independencia de B
    if (B) { S }
}
\end{lstlisting}

\noindent Estos casos no son anecdóticos: son precisamente los que obligan a
definir precondiciones explícitas y verificables en lugar de aplicar la
transformación de forma indiscriminada.

\section{Familia de problemas y alcance}
\label{sec:alcance}

La unidad de transformación es un par de \lstinline|if| anidados. Una cadena de
$N$ \lstinline|if| puramente anidados (cada nivel un único \lstinline|if|, sin
\lstinline|else| ni sentencias intercaladas) se reduce aplicando la regla de
combinación de forma repetida, un par cada vez, hasta que una precondición lo
impida; por eso una cadena puede combinarse total o parcialmente: si un nivel
intermedio incumple P1--P5, la combinación se detiene en ese punto. Quedan
fuera del alcance, como extensiones naturales (\autoref{cap:conclusiones}), las
variantes que no son anidamiento puro: cadenas con \lstinline|else if|,
secuencias de \lstinline|if| hermanos (no anidados) o combinaciones que exigen
reordenar el flujo.

La \autoref{tab:alcance} resume qué entra y qué queda fuera del alcance
inicial. Las construcciones excluidas no son irrelevantes; simplemente son
aquellas cuya seguridad no puede establecerse con garantías dentro del alcance
del prototipo.

\begin{table}[ht]
\centering
\caption{Delimitación del alcance del patrón objetivo.}
\label{tab:alcance}
\begin{tabular}{@{}p{6.1cm}p{6.1cm}@{}}
\toprule
\textbf{Dentro del alcance} & \textbf{Fuera del alcance inicial} \\
\midrule
Un \lstinline|if| que contiene exactamente otro \lstinline|if| (anidamiento simple)
  & Ramas \lstinline|else| o \lstinline|else if| en cualquiera de los dos \lstinline|if| \\
Ausencia de sentencias adicionales en el bloque externo
  & Sentencias antes o después del \lstinline|if| interno \\
Condiciones \lstinline|A| y \lstinline|B| sin efectos colaterales
  & Efectos colaterales o asignaciones en las condiciones \\
Cuerpo \lstinline|S| arbitrario
  & Expresiones lambda, \emph{streams}, concurrencia \\
Combinación mediante \lstinline|&&| con preservación del cortocircuito
  & Construcciones cuya equivalencia no se garantiza sintácticamente \\
\bottomrule
\end{tabular}
\end{table}